% \newdualentry{label}{acronym}{long form}{description}

\newdualentry{dac}{DAC}{Dynamic Airspace Configuration}{%
    è un concetto sviluppato da \citet{eurocontrol2025} e dal programma \gls{sesar} per rendere la gestione dello spazio aereo più flessibile e adattabile. Consiste nella riorganizzazione dinamica dei settori di controllo del traffico aereo in risposta alle variazioni della domanda di traffico, alle condizioni meteorologiche e ad altri vincoli operativi
}

\newdualentry{atm}{ATM}{Air Traffic Management}{%
    è un insieme di funzioni svolte a terra e in volo (servizi del traffico aereo, gestione dello spazio aereo e gestione dei flussi di traffico aereo) finalizzate a garantire il movimento sicuro ed efficiente degli aeromobili durante tutte le fasi delle operazioni
}

\newdualentry{atc}{ATC}{Air Traffic Control}{%
    è un servizio fornito allo scopo di prevenire collisioni tra aeromobili, nonché tra aeromobili e ostacoli nell'area di manovra, e di garantire un flusso di traffico aereo spedito e ordinato
}


\newdualentry{ansp}{ANSP}{Air Navigation Service Provider}{%
    è un organismo, pubblico o privato, responsabile dell’erogazione dei servizi di traffico aereo, gestione dello spazio aereo, comunicazioni, navigazione e sorveglianza, secondo quanto previsto dalle normative e dalle autorità nazionali competenti
}

\newdualentry{acc}{ACC}{Area Control Centre}{%
     è l ’unità operativa responsabile della fornitura del servizio di controllo del traffico aereo agli aeromobili che operano all'interno di uno spazio aereo controllato, tipicamente costituito da una o più \gls{cta}
}

\newdualentry{asm}{ASM}{Airspace Management}{%
    è un processo di progettazione, suddivisione e assegnazione flessibile dello spazio aereo fra usi civili e militari, articolato su tre livelli: strategico (L1), pre-tattico (L2) e tattico (L3)
}

\newdualentry{atfcm}{ATFCM}{Air Traffic Flow \& Capacity Management}{%
    è un'insieme di misure strategiche, pre-tattiche e tattiche tramite cui il Network Manager bilancia in modo continuo la domanda di traffico con la capacità del sistema \gls{atm}
}

\newdualentry{dmean}{DMEAN}{Dynamic Management of the European Airspace Network}{%
    è un è programma di \citet{eurocontrol2025} volto a ottimizzare l’utilizzo dello spazio aereo coordinando in tempo quasi reale i processi di \gls{asm} e \gls{atfcm}
}

\newdualentry{sesar}{SESAR}{Single European Sky ATM Research}{%
    è un programma dell’UE gestito dalla SESAR Joint Undertaking, partenariato tra Commissione, \citet{eurocontrol2025} e industria, che sviluppa soluzioni per modernizzare la gestione del traffico aereo europeo
}

\newdualentry{aixm}{AIXM}{Aeronautical Information Exchange Model}{%
è lo standard internazionale per la codifica e lo scambio digitale di informazioni aeronautiche
}

\newacronym{adrr}{ADRR}{Aviation Data Repository for Research}
\newacronym{adr}{ADR}{Airspace Data Repository}
\newacronym{fl}{FL}{FLight Level}
\newacronym{cta}{CTA}{Control Area}
\newacronym{fua}{FUA}{Flexible Use of Airspace}
\newacronym{fir}{FIR}{Flight Information Region}
\newacronym{aua}{AUA}{ATC Unit Airspace}
\newacronym{airac}{AIRAC}{Aeronautical Information Regulation and Control}